\section{Research target, source boundary, and conventions}
\label{sec:scope}

\subsection{What is continued here}

The source for this investigation is the ProveIt manuscript and its six
current incoming reports at commit
\begin{center}\small\repoSHA.\end{center}
The source was inspected on 11 October 2026. The canonical manuscript
already develops polylogarithmic identities, harmonic generating
functions, Gamma correlations, and several finite-part calculi
\cite{repo}. The incoming Bernoulli--Tornheim--digamma report proves
optimal-prefix rank for the unequal even tail-order pairs $(2,q)$,
$(4,q)$, and $(6,q)$, and asks for every reconstruction row in closed
form \cite[research question R2]{incoming-bernoulli}. The first rational
row is already known. This paper supplies the full inverse for $(2,q)$
and turns it into a family of evaluated ordinary series.

The six incoming archives were also inspected to distinguish new work
from statements already supplied there. In particular, the general
ordered-partition expansion of a product of tails, the symmetric cubic
descent, the finite differential-polynomial integration-by-parts
algorithm, and the first Gauss-determinant resonance are prior parts of
this programme \cite{incoming-cubic,incoming-differential,
incoming-determinantal,incoming-tail,incoming-multishift}.
They are not relabelled as new discoveries in this article.

The word ``new'' below refers to the advancement beyond these inspected
sources. A search of the relevant primary literature was used to identify
classical ingredients; it does not establish global historical priority.
The Bernoulli Hankel and orthogonal-polynomial background is classical
\cite{dilcher-jiu,dlmf-op}. We provide the particular algebraic proofs
needed here because the exceptional zeroth Bernoulli moment and the
tail matrix's omitted binomial terms affect the answer.

\subsection{Three different meanings of an evaluation}

We distinguish an ordinary convergent sum, a meromorphically continued
value, and a constant term of a specified cutoff. Where two of them
coincide, that coincidence is proved. In the centered two-tail problem,
the identities $\zeta(-2j,\tfrac12)=0$ make the requisite polynomial
subtractions have zero continued sum. That special cancellation does
not hold after an arbitrary change of centering.

The numerical checks in the accompanying files are diagnostics, not
interval certificates. An identity in a polynomial ring over $\Q$ is
checked exactly. A theorem about an infinite family or a continuum of
parameters is justified by the written proof. No result in the paper
uses a failed integer-relation search as an independence argument.

\subsection{Notation}

The Gamma function is $\Gamma$, the digamma function is
$\psi=\Gamma'/\Gamma$, and $\psi^{(m)}$ is its $m$th derivative.
For $p>1$ and $n\geq0$, write
\begin{equation}\label{eq:tails}
 H_n^{(p)}=\sum_{j=1}^n j^{-p},\qquad
 A_p(n)=\zeta(p)-H_n^{(p)}=\zeta(p,n+1),\qquad
 x_n=n+\tfrac12.
\end{equation}
For integer $p\geq2$,
\begin{equation}\label{eq:tail-polygamma}
 A_p(n)=\frac{(-1)^p}{(p-1)!}\psi^{(p-1)}(n+1).
\end{equation}
The Hurwitz-zeta and polygamma conventions agree with
\cite{dlmf-hurwitz,dlmf-polygamma}.
For $a>0$, the generalized Stieltjes constants are normalized by
\begin{equation}\label{eq:stieltjes-convention}
 \zeta(1+u,a)=\frac1u+
   \sum_{j=0}^{\infty}\frac{(-1)^j}{j!}\gamma_j(a)u^j.
\end{equation}
Thus $\gamma_0(a)=-\psi(a)$ and $\gamma_j=\gamma_j(1)$.

The Bernoulli convention is
\begin{equation}\label{eq:bernoulli-convention}
 \frac{te^{xt}}{e^t-1}=\sum_{j=0}^{\infty}B_j(x)\frac{t^j}{j!},
 \qquad B_j=B_j(0),\qquad B_1=-\tfrac12.
\end{equation}
The notation $(a)_j$ means a rising factorial. A falling factorial is
$a^{\underline j}=a(a-1)\cdots(a-j+1)$. Empty products are one.

For $|z|<1$, our depth-two polylogarithm is
\begin{equation}\label{eq:depth-two-convention}
 \Li_{p,q}(z)=\sum_{n>m\geq1}\frac{z^n}{n^p m^q}.
\end{equation}
The color belongs to the larger index. This convention matters in the
generating-function identities. At convergent endpoints we write
$\zeta_{\mathrm{MZ}}(p,q)=\sum_{n>m\geq1}n^{-p}m^{-q}$ for a
multiple zeta value; this differs from the Hurwitz function $\zeta(s,a)$.
The spectral double series $Z(s,t)$ in Section~\ref{mellin:section}
uses the same strict ordering on its stated convergence domain.
Auxiliary coefficients such as $C_j$ and $R_j$ are defined locally in
each family and do not identify constants from different sections.
All functions on the open unit disc are
defined by their series; real boundary limits are taken with $z\uparrow1$.
Complex powers of positive real numbers use the real logarithm. Additional
complex-parameter domains are stated locally where they are used.
